\documentclass[10pt,a4paper]{amsart}
\usepackage{amsmath,amssymb,amsthm,graphicx}
\usepackage{verbatim}
\newenvironment{reference}{\comment}{\endcomment}
%\newenvironment{reference}{}{}
\newenvironment{slogan}{\comment}{\endcomment}
\newenvironment{history}{\comment}{\endcomment}

% For commutative diagrams we use Xy-pic
\usepackage[all]{xy}

% We use 2cell for 2-commutative diagrams.
\xyoption{2cell}
% Theorem environments.
%
\theoremstyle{plain}
\newtheorem{theorem}[subsection]{Theorem}
\newtheorem{proposition}[subsection]{Proposition}
\newtheorem{lemma}[subsection]{Lemma}

\theoremstyle{definition}
\newtheorem{definition}[subsection]{Definition}
\newtheorem{example}[subsection]{Example}
\newtheorem{exercise}[subsection]{Exercise}
\newtheorem{situation}[subsection]{Situation}

\theoremstyle{remark}
\newtheorem{remark}[subsection]{Remark}
\newtheorem{remarks}[subsection]{Remarks}

\numberwithin{equation}{subsection}

% Macros
%
\def\lim{\mathop{\mathrm{lim}}\nolimits}
\def\colim{\mathop{\mathrm{colim}}\nolimits}
\def\Spec{\mathop{\mathrm{Spec}}}
\def\Hom{\mathop{\mathrm{Hom}}\nolimits}
\def\Ext{\mathop{\mathrm{Ext}}\nolimits}
\def\SheafHom{\mathop{\mathcal{H}\!\mathit{om}}\nolimits}
\def\SheafExt{\mathop{\mathcal{E}\!\mathit{xt}}\nolimits}
\def\Sch{\mathit{Sch}}
\def\Mor{\mathop{\mathrm{Mor}}\nolimits}
\def\Ob{\mathop{\mathrm{Ob}}\nolimits}
\def\Sh{\mathop{\mathit{Sh}}\nolimits}
\def\NL{\mathop{N\!L}\nolimits}
\def\CH{\mathop{\mathrm{CH}}\nolimits}
\def\proetale{{pro\text{-}\acute{e}tale}}
\def\etale{{\acute{e}tale}}
\def\QCoh{\mathit{QCoh}}
\def\Ker{\mathop{\mathrm{Ker}}}
\def\Im{\mathop{\mathrm{Im}}}
\def\Coker{\mathop{\mathrm{Coker}}}
\def\Coim{\mathop{\mathrm{Coim}}}

% Boxtimes
%
\DeclareMathSymbol{\boxtimes}{\mathbin}{AMSa}{"02}

%
% Macros for moduli stacks/spaces
%
\def\QCohstack{\mathcal{QC}\!\mathit{oh}}
\def\Cohstack{\mathcal{C}\!\mathit{oh}}
\def\Spacesstack{\mathcal{S}\!\mathit{paces}}
\def\Quotfunctor{\mathrm{Quot}}
\def\Hilbfunctor{\mathrm{Hilb}}
\def\Curvesstack{\mathcal{C}\!\mathit{urves}}
\def\Polarizedstack{\mathcal{P}\!\mathit{olarized}}
\def\Complexesstack{\mathcal{C}\!\mathit{omplexes}}
% \Pic is the operator that assigns to X its picard group, usage \Pic(X)
% \Picardstack_{X/B} denotes the Picard stack of X over B
% \Picardfunctor_{X/B} denotes the Picard functor of X over B
\def\Pic{\mathop{\mathrm{Pic}}\nolimits}
\def\Picardstack{\mathcal{P}\!\mathit{ic}}
\def\Picardfunctor{\mathrm{Pic}}
\def\Deformationcategory{\mathcal{D}\!\mathit{ef}}
\usepackage{cleveref}
\crefname{originalproposition}{original proposition}{original propositions}
\Crefname{originalproposition}{Original proposition}{Original propositions}
\makeatletter
\newlabel{algebra-proposition-dimension}{{00KQ}{source}{}{sourcepointer-1}{}}
\newlabel{algebra-proposition-dimension@cref}{{[originalproposition][0][]00KQ}{[1][0][]source}{}{}{}}
\makeatother
\crefname{originalreference}{original reference}{original references}
\Crefname{originalreference}{Original reference}{Original references}
\makeatletter
\newlabel{lemma-locally-Noetherian-quasi-separated}{{01OY}{source}{}{sourcepointer-2}{}}
\newlabel{lemma-locally-Noetherian-quasi-separated@cref}{{[originalreference][0][]01OY}{[1][0][]source}{}{}{}}
\makeatother
\crefname{originallemma}{original lemma}{original lemmas}
\Crefname{originallemma}{Original lemma}{Original lemmas}
\makeatletter
\newlabel{schemes-lemma-specialize-points}{{01J7}{source}{}{sourcepointer-3}{}}
\newlabel{schemes-lemma-specialize-points@cref}{{[originallemma][0][]01J7}{[1][0][]source}{}{}{}}
\makeatother
\title{Every infinite subset of a Noetherian scheme contains an infinite subset without nontrivial specializations}
\author{The Stacks Project Authors}
\date{stacks-project:properties.tex:local-snapshot-20261012; GNU FDL 1.2 or later}
\begin{document}
\maketitle
\noindent\textit{Editorial note.} All delivered mathematical text is the author's own, copied from the pinned Stacks Project chapter file properties.tex; no statement or proof step was rewritten, reordered or removed.\par\medskip
\noindent\textit{Editorial note.} The counted proof is the author's complete argument: the separation of the points with no nontrivial specialization into T, the inductive definition of the strata with the dimension bound showing that only finitely many are nonempty, and the Noetherianity contradiction ruling out points outside the strata.\par\medskip
\noindent\textit{Editorial note.} The label lemma-thin-infinite-sequence names the counted statement, and the proof environment immediately following it in the source is delivered verbatim as its proof.\par\medskip
\noindent\textit{Editorial note.} The counted statement is carried by the author's own lemma environment in the chapter source.\par\medskip
\noindent\textit{Editorial note.} Pre-open gate: python operations/agent-authoring/tools/precard\_gate.py --source properties.tex --label lemma-thin-infinite-sequence --start 843 --end 893 --local-id 28032 returned RESULT=OK (L28032 free, occupancy FREE, public interval clear). Candidate taken from sources/stacks-properties/candidates-halfA.json.\par\medskip
\noindent\textit{Editorial note.} The statement and proof are copied from the pinned chapter file; the cited inputs are kept as external pointers with their canonical Stacks tags.\par\medskip
\noindent\textit{Allowed standard prerequisites (editorial).}\begin{itemize}
\item Finiteness of the dimension of the local rings of a Noetherian scheme. Algebra, the proposition that the dimension of the local ring of a Noetherian scheme at a point is finite. Supplies the finiteness of the maximum dimension d of the local rings at the points of the finite set T\_0.
\item Realization of the points specializing to a point inside its local ring. Schemes, the lemma that the points specializing to a given point are exactly the points of the spectrum of the local ring at that point. Supplies the identification of the points of a chain of specializations ending in T\_0 as points of the spectrum of the local ring there, which bounds the length of the chain.
\item Noetherianness of the underlying topological space. Same chapter, the lemma that a Noetherian scheme has a Noetherian underlying topological space. Supplies the impossibility of an infinite descending chain of specializations in the complement of the finite union.
\end{itemize}\medskip
\ifdefined\hypertarget\hypertarget{sourcepointer-1}{}\fi
\noindent\textit{Source reference.} Original proposition 00KQ: sources/stacks-chow/tags.txt, source lines 760--760. Prerequisite pointer invoked at proof line(s) 868; it names a standard fact of a sibling chapter that the counted argument uses, and every step of the counted proof built on it is retained verbatim in the delivered body.\par\medskip
\ifdefined\hypertarget\hypertarget{sourcepointer-2}{}\fi
\noindent\textit{Source reference.} Original reference 01OY: sources/stacks-properties/pinned-original-source.tex, source lines 698--698. Prerequisite pointer invoked at proof line(s) 892; it names a standard fact of the same chapter that the counted argument uses, and every step of the counted proof built on it is retained verbatim in the delivered body.\par\medskip
\ifdefined\hypertarget\hypertarget{sourcepointer-3}{}\fi
\noindent\textit{Source reference.} Original lemma 01J7: sources/stacks-chow/tags.txt, source lines 2004--2004. Prerequisite pointer invoked at proof line(s) 876; it names a standard fact of a sibling chapter that the counted argument uses, and every step of the counted proof built on it is retained verbatim in the delivered body.\par\medskip
\begin{lemma}
\label{lemma-thin-infinite-sequence}
Let $S$ be a Noetherian scheme. Let $T \subset S$ be an infinite subset.
Then there exists an infinite subset $T' \subset T$
such that there are no nontrivial specializations among the points $T'$.
\end{lemma}

\begin{proof}
Let $T_0 \subset T$ be the set of $t \in T$ which do not specialize
to another point of $T$. If $T_0$ is infinite, then $T' = T_0$ works.
Hence we may and do assume $T_0$ is finite.
Inductively, for $i > 0$, consider the set $T_i \subset T$
of $t \in T$ such that
\begin{enumerate}
\item $t \not \in T_{i - 1} \cup T_{i - 2} \cup \ldots \cup T_0$,
\item there exist a nontrivial specialization $t \leadsto t'$ with
$t' \in T_{i - 1}$, and
\item for any nontrivial specialization
$t \leadsto t'$ with $t' \in T$ we have
$t' \in T_{i - 1} \cup T_{i - 2} \cup \ldots \cup T_0$.
\end{enumerate}
Again, if $T_i$ is infinite, then $T' = T_i$ works.
Let $d$ be the maximum of the dimensions of the local rings
$\mathcal{O}_{S, t}$ for $t \in T_0$; then $d$ is an integer
because $T_0$ is finite and the dimensions of the local rings
are finite by Algebra, Proposition \ref{algebra-proposition-dimension}.
Then $T_i = \emptyset$ for $i > d$.
Namely, if $t \in T_i$ then we can find a sequence
of nontrivial specializations
$t = t_i \leadsto t_{i - 1} \leadsto \ldots \leadsto t_0$
with $t_0 \in T_0$. As
the points $t = t_i, t_{i - 1}, \ldots, t_0$ are in
$\Spec(\mathcal{O}_{S, t_0})$
(Schemes, Lemma \ref{schemes-lemma-specialize-points}),
we see that $i \leq d$.
Thus $\bigcup T_i = T_d \cup \ldots \cup T_0$ is a finite subset of $T$.

\medskip\noindent
Suppose $t \in T$ is not in $\bigcup T_i$. Then there must be a specialization
$t \leadsto t'$ with $t' \in T$ and $t' \not \in \bigcup T_i$. (Namely, if
every specialization of $t$ is in the finite set $T_d \cup \ldots \cup T_0$,
then there is a maximum $i$ such that there is some specialization
$t \leadsto t'$ with $t' \in T_i$ and then $t \in T_{i + 1}$ by construction.)
Hence we get an infinite sequence
$$
t \leadsto t' \leadsto t'' \leadsto \ldots
$$
of nontrivial specializations between points of $T \setminus \bigcup T_i$.
This is impossible because the underlying topological space of $S$
is Noetherian by Lemma \ref{lemma-locally-Noetherian-quasi-separated}.
\end{proof}

\end{document}
